\chapter{How the machine sees data}
\label{ch:02}

Chapter~1 followed a program from text to running code. This chapter asks
the companion question: once it is running, \emph{what does the machine
think the data is}? The honest answer is short and unsettling: nothing.
A memory cell holds bytes; a register holds bits. Types --- signed or
unsigned, integer or pointer or float --- are agreements between your
source code and the generated instructions, and the CPU keeps no record
of them. Most of this chapter is spent making that concrete: we decode
the same four bytes three ways, measure what struct layout really costs,
and run an alignment experiment whose result surprised us --- the numbers
are real, and they disagree with the folklore. That disagreement is the
lesson.

\begin{wins}[What you will be able to do after this chapter]
You will be able to take a raw hex dump --- a line of bytes with no
labels --- and reconstruct what the program that wrote it meant: which
bytes form which integers (including negatives, via two's complement),
which byte order the machine used, where the compiler inserted padding
inside a struct, and how to read a 32-bit pattern as a floating-point
number. You will also be able to say what a pointer really is --- an
address with a type --- and why an array name is almost, but not quite,
the same thing.
\end{wins}

\begin{rulebox}[Before you start]
One thing carries over; it is refreshed here.
\emph{Binary and hex} (Book 5, ch01--02): a byte is eight bits; hex is
just base-16 shorthand for four bits at a time, so \texttt{0x2F} is the
byte \texttt{0010\,1111}. You do not need to convert fluently --- the
chapter rebuilds the arithmetic as it goes --- only to believe that a
hex dump is a picture of memory, not a code to be cracked.
\emph{A compiled program} (ch01 of this book): you saw source become
machine code. Here we look at the \emph{data} that code moves around.
\end{rulebox}

\index{endianness}\index{alignment}\index{two's complement}%
\index{IEEE 754}\index{floating point}\index{struct padding}%
\index{pointer}\index{array}\index{LP64}\index{type punning}

\section{Bytes, words and addresses}
\label{sec:bytes02}

\index{word}\index{address}
A \emph{byte} is the smallest thing a modern machine addresses: 8 bits,
one address. A \emph{word} is the machine's working size --- 32 bits on
RV32, 64 on the laptop you are probably using. When a program says
``the value at address \texttt{0x2000}'', it means ``the byte at
\texttt{0x2000}''; a 4-byte \texttt{int} stored there occupies
\texttt{0x2000} through \texttt{0x2003} and belongs to whichever
interpretation reads it. The machine's data types table is short enough
to memorise:

\begin{figure}[htbp]
\centering
\begin{center}
\begin{tabularx}{0.9\linewidth}{lXXr}
\toprule
\textbf{C type} & \textbf{LP64 (Linux, macOS arm64)}
                & \textbf{ILP32 (RV32, Win32)} & \textbf{bits} \\
\midrule
\texttt{char}        & 1 & 1 & 8 \\
\texttt{short}       & 2 & 2 & 16 \\
\texttt{int}         & 4 & 4 & 32 \\
\texttt{long}        & 8 & 4 & 64 / 32 \\
\texttt{long long}   & 8 & 8 & 64 \\
\texttt{float}       & 4 & 4 & 32 \\
\texttt{double}      & 8 & 8 & 64 \\
\texttt{void *}      & 8 & 4 & 64 / 32 \\
\bottomrule
\end{tabularx}
\end{center}

\caption{C type sizes: LP64 (what this book's dev box runs) vs ILP32
(the RV32 model the capstone targets). \texttt{long} and pointers are
the two that move.}
\label{fig:typetable02}
\end{figure}

\begin{tripwire}[\texttt{long} is a trap on every platform]
On LP64 \texttt{long} is 8 bytes; on ILP32 --- including RV32, the ISA
your capstone CPU runs --- it is 4. Code that assumes \texttt{sizeof(long)
== 8} breaks silently when cross-compiled. Use \texttt{int32\_t} /
\texttt{uint64\_t} from \texttt{<stdint.h>} whenever the size matters,
and it always matters in this book.
\end{tripwire}

\section{Two's complement, in one page}
\label{sec:twos02}

\index{sign bit}\index{overflow}
Book~5 ch8 derived this properly; here is the working version. An
$n$-bit signed integer stores $-2^{n-1} + \sum b_i 2^i$ --- the top bit
contributes \emph{minus} its weight. That single negative weight buys
three things: one zero, one adder for both addition and subtraction,
and a clean wrap from most-positive to most-negative. Figure
\ref{fig:twos02} shows the 8-bit case on a line: $127 + 1$ does not
saturate, it jumps to $-128$.

\begin{figure}[htbp]
\centering
\begin{tikzpicture}[font=\scriptsize, >=latex]
% two's complement wrap-around number line, int8
\draw[->, thick] (0.4,2.6) -- (9.9,2.6);
\foreach \x/\v/\u in {
  0.6/{0000 0000}/{0},
  2.0/{0000 0001}/{1},
  3.4/{0111 1111}/{127},
  4.8/{1000 0000}/{-128},
  6.2/{1000 0001}/{-127},
  8.4/{1111 1110}/{-2},
  9.4/{1111 1111}/{-1}}{
  \draw[thick] (\x,2.5) -- (\x,2.7);
  \node[font=\tiny, above] at (\x,2.75) {\v};
  \node[font=\tiny, below] at (\x,2.45) {\u};
}
\draw[->, very thick, accent] (3.4,3.4) to[bend left=45] (4.8,3.4);
\node[font=\tiny, text=accent] at (4.1,4.1) {$127+1 \to -128$};
\node[align=left, text=slate, text width=10cm] at (5.0,1.6)
  {Two's complement is one hardware story, not two: adding and
   subtracting are the same circuit, and the ``wrap'' at the ends is
   why signed overflow is a cliff, not a curve.};
\end{tikzpicture}

\caption{Two's complement as a number line: the wrap at the ends is
hardware fact, not an error message.}
\label{fig:twos02}
\end{figure}

\begin{principle}
The CPU does arithmetic on \emph{bit patterns}, not on numbers. Whether
\texttt{0x80000000} is $-2^{31}$ or $+2^{31}$ is decided entirely by the
instruction that consumes it --- \texttt{slt} vs \texttt{sltu} in RV32I,
same bits, two opinions.
\end{principle}

\section{Endianness: which byte goes first}
\label{sec:endian02}

\index{little-endian}\index{big-endian}
A 4-byte value has to go into 4 one-byte slots, and there are two
conventions. Little-endian puts the least-significant byte at the
lowest address; big-endian puts the most-significant first. RV32 is
little-endian, x86 is little-endian, ARM defaults to little --- the
network protocols you will meet in other books are big-endian, which is
why the conversion functions are called \texttt{htonl}/\texttt{ntohl}.

\begin{example}[Real bytes on real hardware]
\label{ex:bytes02}
\texttt{code/ch02\_bytes.c} stores \texttt{0x0A0B0C0D} in a
\texttt{uint32\_t} and prints the four memory bytes in address order.
Real output on the dev box (arm64, LP64):

\begin{lstlisting}[basicstyle=\ttfamily\scriptsize]
bytes (LE order): 0d 0c 0b 0a
as int32 : 168496141
as float : 6.69487e-33
0x3FC00000 as float: 1.5
\end{lstlisting}

The least-significant byte \texttt{0x0d} lands first: little-endian,
confirmed by experiment rather than documentation. The same four bytes
then get read as a signed int (\texttt{168496141} --- same value, the
top bit happened to be zero) and as a float (nonsense:
\texttt{6.69e-33} --- the exponent field is tiny). Three different
worlds from one bit pattern; the next section formalises the third.
\end{example}

\begin{figure}[htbp]
\centering
\begin{tikzpicture}[font=\scriptsize, >=latex]
% 0x0A0B0C0D little vs big endian
\node[align=left, font=\footnotesize\bfseries] at (1.7,4.6)
  {little-endian (RV32, x86, ARM)};
\node[align=left, font=\footnotesize\bfseries] at (7.0,4.6)
  {big-endian (network order)};
\foreach \i/\v in {0/0D,1/0C,2/0B,3/0A}{
  \draw[thick, fill=accentsoft] (0.4+\i*0.8,3.8) rectangle (1.2+\i*0.8,4.4);
  \node at (0.8+\i*0.8,4.1) {\v};
  \node[font=\tiny, text=slate] at (0.8+\i*0.8,3.6) {+\i};
}
\foreach \i/\v in {0/0A,1/0B,2/0C,3/0D}{
  \draw[thick, fill=graybg] (5.7+\i*0.8,3.8) rectangle (6.5+\i*0.8,4.4);
  \node at (6.1+\i*0.8,4.1) {\v};
  \node[font=\tiny, text=slate] at (6.1+\i*0.8,3.6) {+\i};
}
\node[align=left, font=\tiny, text width=4.4cm] at (1.7,3.0)
  {LSB lands at the lowest address --- same value, reversed paper trail};
\node[align=left, font=\tiny, text width=4.4cm] at (7.0,3.0)
  {MSB first --- reads left-to-right like the number you wrote};
\node[align=left, text=slate, text width=10cm] at (5.0,2.2)
  {Real bytes from \ttfamily ch02\_bytes\rmfamily: \ttfamily 0d 0c 0b 0a
   \rmfamily --- your dev box writes little-endian, and so does RV32.};
\end{tikzpicture}

\caption{\texttt{0x0A0B0C0D} in memory, both conventions. RV32, x86 and
ARM all pick the left column.}
\label{fig:endian02}
\end{figure}

\begin{checkpoint}
You receive the four bytes \texttt{78 56 34 12} (in address order) over
a little-endian memory-mapped register. What 32-bit value did the
sender write?
\end{checkpoint}
\begin{solution}
\texttt{0x12345678} --- little-endian means the first byte is the
least significant, so read the sequence backwards as hex pairs.
\end{solution}

\section{Alignment and struct layout}
\label{sec:align02}

\index{padding}\index{offsetof@\texttt{offsetof}}\index{aligned access}
Hardware prefers a $n$-byte value at an address divisible by $n$ ---
an \emph{aligned} access. Compilers honour this by inserting
\emph{padding}: invisible bytes between struct fields so every field
sits at its natural offset. The rule chain: each field's offset must be
a multiple of its alignment; the struct's size is rounded up to the
largest member's alignment; arrays of the struct need that tail padding
to keep every element aligned.

\begin{example}[The same three fields, two prices]
\label{ex:pad02}
\texttt{code/ch02\_pad.c} declares \texttt{struct A\{char;double;int;\}}
and \texttt{struct B\{double;int;char;\}} --- identical fields, different
order. Real output:

\begin{lstlisting}[basicstyle=\ttfamily\scriptsize]
A: sizeof=24  c@0 d@8 i@16
B: sizeof=16  d@0 i@8 c@12
\end{lstlisting}

In \texttt{A}, \texttt{char c} occupies byte 0, then \texttt{double d}
must start at an 8-aligned address: bytes 1--7 are padding. \texttt{int i}
lands at 16, and four tail bytes pad the total to 24. Reordering the
same three fields into \texttt{B} costs 16 bytes --- a third less
memory, for free, just by sorting fields largest-first.
\end{example}

\begin{figure}[htbp]
\centering
\begin{tikzpicture}[font=\scriptsize, >=latex]
% struct A vs B padding map (real offsetof output)
\node[align=left, font=\footnotesize\bfseries] at (2.0,4.6)
  {\ttfamily struct A\{char c; double d; int i;\}};
\node[align=left, font=\footnotesize\bfseries] at (8.0,4.6)
  {\ttfamily struct B\{double d; int i; char c;\}};
% A layout: c@0 pad1-7 d@8-15 i@16-19 pad20-23
\draw[thick, fill=accentsoft] (0.4,3.8) rectangle (0.9,4.4);
\node[font=\tiny] at (0.65,4.1) {c};
\draw[thick, fill=graybg]
  (0.9,3.8) rectangle (2.4,4.4);
\node[font=\tiny] at (1.65,4.1) {pad 7 B};
\draw[thick, fill=goodbg] (2.4,3.8) rectangle (5.4,4.4);
\node[font=\tiny] at (3.9,4.1) {d (8 B, needs 8-aligned)};
\draw[thick, fill=accentsoft] (5.4,3.8) rectangle (6.9,4.4);
\node[font=\tiny] at (6.15,4.1) {i (4 B)};
\draw[thick, fill=graybg]
  (6.9,3.8) rectangle (7.9,4.4);
\node[font=\tiny] at (7.4,4.1) {pad};
\node[font=\tiny, text=slate] at (4.15,3.5)
  {\ttfamily sizeof=24 --- real \ttfamily offsetof\rmfamily output};
% B layout: d@0-7 i@8-11 c@12 pad13-15
\draw[thick, fill=goodbg] (6.2,2.6) rectangle (9.2,3.2);
\node[font=\tiny] at (7.7,2.9) {d (8 B)};
\draw[thick, fill=accentsoft] (9.2,2.6) rectangle (10.4,3.2);
\node[font=\tiny] at (9.8,2.9) {i (4 B)};
\draw[thick, fill=accentsoft] (10.4,2.6) rectangle (10.8,3.2);
\node[font=\tiny] at (10.6,2.9) {c};
\node[font=\tiny, text=slate] at (8.5,2.3)
  {\ttfamily sizeof=16 --- same fields, 8 bytes cheaper};
\node[align=left, text=slate, text width=10.4cm] at (5.2,1.5)
  {The compiler pads to each field's alignment and the struct's worst
   member; reordering fields is free bytes. Measured, not theorised.};
\end{tikzpicture}

\caption{Padding maps of \texttt{A} vs \texttt{B}: measured
\texttt{offsetof} values, not a sketch.}
\label{fig:padding02}
\end{figure}

\subsection{Measurement: the alignment tax that wasn't}
\label{sec:measure02}

Folklore says misaligned access is slow: the CPU splits the load across
two cache lines and pays double. \texttt{code/ch02\_align.c} tests it
directly --- sum a 4\,M-element \texttt{uint32} array once through an
aligned pointer, once through the same pointer offset by one byte
(every load misaligned). Real timings, Apple M-series, 32 passes:

\begin{lstlisting}[basicstyle=\ttfamily\scriptsize]
offset 0: 151.2 ms  (s=2260630400663552)
offset 1: 110.4 ms  (s=4521260264456192)
\end{lstlisting}

The misaligned version is \emph{faster} --- repeatably, in both loop
orders, by $\sim$25\%. Two honest conclusions. First, the expected tax
is absent: Apple Silicon performs unaligned loads at full speed, as do
most modern cores. Second --- the reason this section exists --- a
consistent 25\% delta in the \emph{wrong direction} means the benchmark
is measuring something else (loop layout, prefetcher behaviour, uarch
quirks), not alignment. When the data contradicts the folklore, the
data wins and the model gets revised. That is the book's method in
miniature: predict, measure, reconcile.

\begin{tripwire}[``Misaligned is slow'' is hardware folklore]
It was true on 1980s CPUs and still holds on some embedded cores ---
and it is false on the chip in your laptop. Rules about performance
expire; measurements do not. Re-test before you re-tell.
\end{tripwire}

\section{IEEE-754 in one breath}
\label{sec:float02}

\index{mantissa}\index{exponent}\index{biased exponent}
A \texttt{float} is 32 bits partitioned $S\,|\,E\,|\,F$: one sign bit,
an 8-bit biased exponent, a 23-bit fraction with an implied leading 1.
Value: $(-1)^S \times (1 + F/2^{23}) \times 2^{E-127}$ for
$0 < E < 255$; the extremes $E{=}0$ and $E{=}255$ are reserved for
denormals/zero and inf/NaN. Precision is about 7 decimal digits; the
spacing between neighbours near $2^e$ is $2^{e-23}$, and round-to-nearest
keeps each operation within half that --- the ``$.5$ ulp'' you will hear
cited every time someone says ``floats are exact enough''.

\begin{example}[Decode \texttt{0x40490FDB} by hand, then by code]
\label{ex:float02}
$S{=}0$, exponent field \texttt{0x80} ($=128$),
fraction field \texttt{0x490fdb}. Value
$= 1.5707964 \times 2^{128-127} = 3.1415928$ --- $\pi$ to float
precision. \texttt{code/ch02\_float.c} performs the same decode
bit-by-bit and prints:

\begin{lstlisting}[basicstyle=\ttfamily\scriptsize]
sign=0 exp=128(E-1) frac=0x490fdb
value = 3.141592741
\end{lstlisting}

The last digit differs from $\pi$ because the 23-bit fraction itself is
the approximation --- the stored value \emph{is} the rounding.
\end{example}

\begin{figure}[htbp]
\centering
\begin{tikzpicture}[font=\scriptsize, >=latex]
% IEEE-754 single anatomy for pi = 0x40490FDB
\node[align=left, font=\footnotesize\bfseries] at (4.0,4.6)
  {IEEE-754 single: \texttt{0x40490FDB} ($\pi$, to 7 digits)};
\draw[thick, fill=accentsoft] (0.6,3.8) rectangle (1.6,4.4);
\node at (1.1,4.1) {S};
\node[font=\tiny, text=slate] at (1.1,3.6) {bit 31};
\draw[thick, fill=goodbg] (1.6,3.8) rectangle (4.6,4.4);
\node at (3.1,4.1) {exponent $E$ (8 bits)};
\node[font=\tiny, text=slate] at (3.1,3.6) {bits 30--23, bias 127};
\draw[thick, fill=graybg] (4.6,3.8) rectangle (9.6,4.4);
\node at (7.1,4.1) {fraction $F$ (23 bits)};
\node[font=\tiny, text=slate] at (7.1,3.6) {bits 22--0, hidden leading 1};
\node[align=left, font=\tiny, text width=9.6cm] at (5.1,2.9)
  {value $= (-1)^S \times (1 + F/2^{23}) \times 2^{E-127}$.
   For $\pi$: $S{=}0$, $E{=}128$, $F{=}$\texttt{0x490fdb}
   $\to$ $1.5707964 \times 2^1 = 3.1415927$ --- the decoder in
   \texttt{ch02\_float.c} prints exactly this.};
\node[align=left, text=slate, text width=10cm] at (5.0,1.8)
  {The ``$.5$ ulp'' idea: spacing between neighbours is
   $2^{E-127}\times 2^{-23}$ --- round-to-nearest keeps error under
   half a step, which is why floats lie less than you feared.};
\end{tikzpicture}

\caption{IEEE-754 single precision anatomy, decoded for $\pi$.}
\label{fig:ieee02}
\end{figure}

\begin{checkpoint}
In \texttt{0x3FC00000} (the \texttt{1.5} printed in
Example~\ref{ex:bytes02}), what are $S$, $E$ and the implied mantissa?
\end{checkpoint}
\begin{solution}
$S{=}0$; $E{=}\texttt{0x7F}{=}127$ so the real exponent is $127-127=0$;
$F{=}\texttt{0x400000}$, mantissa $=1 + 0.5 = 1.5$; value
$= 1.5 \times 2^0$.
\end{solution}

\section{Pointers as typed addresses; arrays vs pointers}
\label{sec:ptr02}

\index{pointer arithmetic}
A pointer is an address wearing a type. \texttt{p+1} advances by
\texttt{sizeof(*p)} bytes --- the arithmetic is scaled, which is why
walking an \texttt{int32\_t*} moves 4 bytes per step and a
\texttt{uint8\_t*} moves 1. The type is only a scale factor to the
code generator; the CPU sees an address either way. The difference
between \texttt{int a[4]} and \texttt{int *p} is subtler: the array's
base address is fixed at compile time, while the pointer's value is
itself data --- stored, loaded, changeable. In generated code
\texttt{a[i]} is one load off a known base; \texttt{p[i]} is two: load
\texttt{p}, then load \texttt{p[i]}.

\begin{figure}[htbp]
\centering
\begin{tikzpicture}[font=\scriptsize, >=latex]
% array vs pointer codegen
\draw[thick, fill=accentsoft] (0.3,2.8) rectangle (4.6,4.6);
\node[font=\footnotesize] at (2.45,4.2) {\ttfamily int a[4]};
\node[align=left, font=\tiny, text width=3.9cm] at (2.45,3.4)
  {\texttt{a[i]} $\to$ base $+$ i$\times$4,\\
   one load from a fixed offset ---\\
   the array \emph{is} its address};
\draw[thick, fill=graybg] (5.5,2.8) rectangle (9.8,4.6);
\node[font=\footnotesize] at (7.65,4.2) {\ttfamily int *p};
\node[align=left, font=\tiny, text width=3.9cm] at (7.65,3.4)
  {\texttt{p[i]} $\to$ load p first,\\
   then base $+$ i$\times$4 ---\\
   two memory trips where the array needed one};
\node[align=left, text=slate, text width=10cm] at (5.05,2.0)
  {``Arrays are pointers'' is a half-truth: the array's address is
   compile-time, the pointer's is runtime --- the generated code keeps
   the receipt.};
\end{tikzpicture}

\caption{Array vs pointer in generated code: the array \emph{is} its
address; the pointer must be fetched first.}
\label{fig:ptrarr02}
\end{figure}

\begin{example}[Reading the same bytes three ways, safely]
\label{ex:pun02}
Example~\ref{ex:bytes02} reinterpreted the four bytes of
\texttt{0x0A0B0C0D} as \texttt{int32\_t} (same value, bit-compatible)
and as \texttt{float} ($6.69 \times 10^{-33}$, nonsense) --- both via
\texttt{memcpy}, not pointer casts. That detail matters: type-punning through a
\texttt{float*} cast violates strict aliasing and lets the optimiser
reorder or delete your reads; \texttt{memcpy} is the standards-strict
way to ask ``what do these bytes mean as a different type?'' and
compiles to zero instructions.
\end{example}

\begin{wins}
\begin{itemize}
\item You can say what a byte, a word and an address are --- and that
types are compile-time agreements, not machine facts.
\item You read endianness off real bytes and explain why RV32 chose
little-endian.
\item You predict struct \texttt{sizeof} and \texttt{offsetof} before
running the probe --- and reorder fields to save real bytes.
\item You decode an IEEE-754 bit pattern by hand and with the decoder.
\item You know the alignment-folklore experiment, what it actually
showed, and why ``measure it'' beats ``everyone knows''.
\end{itemize}
\end{wins}

\section{Common mistakes}
\label{sec:mistakes02}

\begin{itemize}
  \item \textbf{``The variable is signed/unsigned/float.''} The
        \emph{bytes} are just bytes; the type lives in which instructions
        touch them. \texttt{0x0A0B0C0D} was an int and a float in the
        same object, at the same address, in the same second.
  \item \textbf{Assuming \texttt{long} is fixed.} It changes size
        between LP64 and ILP32 --- including your RV32 target. Use
        fixed-width types.
  \item \textbf{``Arrays and pointers are the same.''}
        \texttt{sizeof(a)} vs \texttt{sizeof(p)} disagrees, and the
        generated code shows the extra load. Decay happens at call
        boundaries, not in your head.
  \item \textbf{``Misaligned access is slow.''} Measure first ---
        modern cores largely don't care, and our own experiment inverted
        the prediction.
  \item \textbf{Type-punning through casts.} \texttt{*(float*)\&u} is
        strict-aliasing UB; \texttt{memcpy} is free and correct.
  \item \textbf{``Floats can't represent that.''} They represent
        \emph{nearest} values: $0.1$ is not stored, its neighbour is.
        Equality tests belong on integers; floats get tolerances.
\end{itemize}

\section*{Exercises}
\label{sec:exercises02}
\addcontentsline{toc}{section}{Exercises}

\begin{exercise}\label{ex:e02a}
\texttt{uint8\_t} vs \texttt{int8\_t}: which changes when you print the
same bit pattern \texttt{0xFF}, and what does each show?
\end{exercise}
\begin{solution}[hint]
The bits are identical; only the interpretation differs. One prints
255, the other $-1$ --- name which is which.
\end{solution}

\begin{exercise}\label{ex:e02b}
A \texttt{struct} with \texttt{char,int,char} fields reports
\texttt{sizeof=12}. Account for every byte.
\end{exercise}
\begin{solution}[hint]
\texttt{int} needs 4-alignment: c@0, pad 1--3, i@4--7, c@8, tail pad
9--11 so the struct is a multiple of 4.
\end{solution}

\begin{exercise}\label{ex:e02c}
Convert \texttt{0xBF800000} to a decimal float.
\end{exercise}
\begin{solution}[hint]
$S{=}1$, $E{=}127$, $F{=}0$: $-1.0 \times 2^0 = -1.0$.
\end{solution}

\begin{exercise}\label{ex:e02d}
Why does the size of \texttt{"hello"} come out 6 on every
platform, while the size of a \texttt{char*} differs between
data models?
\end{exercise}
\begin{solution}[hint]
The string literal is an array of 6 chars including the NUL; the
pointer is an address, and addresses are the thing that changes size.
\end{solution}

\begin{exercise}\label{ex:m02a}
On a big-endian machine, what would \texttt{ch02\_bytes.c} print in
the first line instead of \texttt{0d 0c 0b 0a}?
\end{exercise}
\begin{solution}[hint]
Reverse the stored order: MSB first ---
\texttt{0a 0b 0c 0d}.
\end{solution}

\begin{exercise}\label{ex:m02b}
The alignment experiment measured misaligned as \emph{faster}. Give
two reasons this could happen that have nothing to do with alignment.
\end{exercise}
\begin{solution}[hint]
Loop ordering/warm-up, code layout, prefetcher behaviour, uarch
quirks, different accumulated work --- name any two; the point is that
a wrong-direction result indicts the benchmark, not the folklore's
opposite.
\end{solution}

\begin{exercise}\label{ex:m02c}
\texttt{double d = 0.1 + 0.2;} prints \texttt{0.30000000000000004}.
Which IEEE-754 fact explains the 4?
\end{exercise}
\begin{solution}[hint]
Neither 0.1 nor 0.2 has a finite binary fraction; both round to a
neighbour, and the sum of the neighbours is one more neighbour ---
the ``4'' is the accumulated rounding surfacing.
\end{solution}

\begin{exercise}\label{ex:m02d}
Your RV32 code does \texttt{*(volatile uint32\_t*)0x10013000} to read
a UART register. Why must the access be exactly 4 bytes, and what
breaks if the compiler splits it into byte loads?
\end{exercise}
\begin{solution}[hint]
MMIO registers are word-sized state machines: reading one byte can
acknowledge the whole register (e.g.\ consume the RX FIFO entry).
Split loads change device behaviour, not just performance.
\end{solution}

\begin{exercise}[hard]\label{ex:h02a}
Write a standards-strict float-bits printer (sign, exponent,
fraction fields) without type-punning UB and without assuming
IEEE-754 --- verify the format macro first.
\end{exercise}
\begin{solution}[hint]
Copy the float's bytes into a \texttt{uint32\_t} with
\texttt{memcpy}, mask the fields, and guard the whole thing with a
\texttt{\_\_FLT\_RADIX\_\_} check --- C never promised IEEE-754; it
just happens to hold everywhere you will ship.
\end{solution}

\begin{exercise}[hard]\label{ex:h02b}
Design a fair alignment benchmark that would actually isolate the
misalignment cost on a machine that hides it. Name three confounders
and how you'd control each.
\end{exercise}
\begin{solution}[hint]
Interleave the two loops (kill ordering effects), keep \texttt{s}
identical (kill work differences), use a dependent-load chain or
volatile store (kill OOO hiding), pin to one core and average many
trials. List any three.
\end{solution}

\section*{Summary}
\label{sec:summary02}
\addcontentsline{toc}{section}{Summary}

\begin{center}
\begin{tabularx}{\linewidth}{lX}
\toprule
\textbf{Idea} & \textbf{One-line version} \\
\midrule
types & compile-time agreements; the CPU sees bytes \\
two's complement & top bit has negative weight; one adder does both \\
endianness & byte order on the wire/in RAM; RV32 is little \\
alignment+padding & fields sit at natural offsets; reorder fields
for free bytes \\
IEEE-754 & $S|E|F$; $\approx$7 digits; neighbours are the truth \\
measurement & folklore expired; our own alignment test inverted \\
\bottomrule
\end{tabularx}
\end{center}

\section*{Further reading}
\label{sec:fr02}
\addcontentsline{toc}{section}{Further reading}

\begin{itemize}
  \item Bryant \& O'Hallaron, \emph{Computer Systems: A Programmer's
        Perspective}, ch.\,2 --- the standard full treatment of every
        section here.
  \item Suzanne J.\ Matthews, Tia Newhall, Kevin C.\ Webb,
        \emph{Dive into Systems}, ch.\,4 --- free online, C-focused
        coverage of representation.
        \url{https://diveintosystems.org/book/C4-Data_rep/}
  \item IEEE-754-2019 standard summary at
        \url{https://en.wikipedia.org/wiki/IEEE_754} --- the format
        table this chapter compresses.
  \item ``What Every Computer Scientist Should Know About Floating-
        Point Arithmetic'' (Goldberg), Sun docs reprint.
  \item[] \href{https://docs.oracle.com/cd/E19957-01/806-3568/ncg_goldberg.html}{docs.oracle.com/cd/E19957-01/806-3568/ncg\_goldberg.html}
  \item RISC-V unprivileged ISA spec, \S2 (base integer model,
        little-endian note) ---
        \url{https://riscv.org/specifications/}
\end{itemize}

\section*{Watch instead}
\label{sec:watch02}
\addcontentsline{toc}{section}{Watch instead}

\begin{rulebox}
\textbf{Watch instead.}
\begin{itemize}
  \item ``Endianness Explained With an Egg'' --- Computerphile, 10:46.
        The boiled-egg walkthrough of Section~\ref{sec:endian02}.
        \url{https://youtu.be/NcaiHcBvDR4}
  \item ``Floating Point Numbers'' --- Computerphile, 9:16. The
        IEEE-754 fields animated; matches Section~\ref{sec:float02}.
        \url{https://youtu.be/PZRI1IfStY0}
\end{itemize}
\end{rulebox}

\chapternotes
